\documentclass[11pt, a4paper]{article}

\usepackage[T1]{fontenc}
\usepackage[utf8]{inputenc}
\usepackage{tgpagella}
\usepackage{tgcursor}
\usepackage{microtype}
\usepackage[spanish, es-noshorthands]{babel}
\addto\captionsspanish{\renewcommand{\tablename}{Tabla}}

\usepackage{amsmath, amssymb, amsfonts, bm}
\usepackage[table, xcdraw]{xcolor}
\usepackage{graphicx}
\usepackage{booktabs}
\usepackage{tabularx}
\usepackage[most]{tcolorbox}
\usepackage[margin=2.5cm]{geometry}
\usepackage{fancyhdr}

\pagestyle{fancy}
\fancyhf{}
\setlength{\headheight}{16pt}
\lhead{\textbf{UCB Sede Tarija} | Ing. Mecatrónica}
\chead{Robótica (IMT-342)}
\rhead{Semana 5 - Sesión 1}
\lfoot{Prof: M.Sc. Ing. Bernardo Quiroga Turdera}
\rfoot{Página \thepage}
\renewcommand{\headrulewidth}{0.4pt}
\renewcommand{\footrulewidth}{0.4pt}

\definecolor{ucbblue}{RGB}{0, 51, 102}

\begin{document}

\begin{center}
    \Large\textbf{\textcolor{ucbblue}{UNIVERSIDAD CATÓLICA BOLIVIANA "SAN PABLO" — SEDE TARIJA}}\\
    \vspace{0.2cm}
    \normalsize Departamento de Ciencias Exactas e Ingenierías | Carrera de Ingeniería Mecatrónica\\
    Asignatura: Robótica Industrial (IMT-342) \\
    Docente: M.Sc. Ing. Bernardo Quiroga Turdera
\end{center}

\vspace{0.4cm}
\begin{tcolorbox}[colback=ucbblue!5!white, colframe=ucbblue, title=\textbf{Hoja de Ejercicios: Cinemática Directa mediante D-H}]
    \centering \large \textbf{Asignación Práctica Independiente}
\end{tcolorbox}
\vspace{0.4cm}

\section*{Ejercicio 1: Evaluación de una fila D-H[cite: 1]}
Dada la convención estándar y los siguientes parámetros para un eslabón $i$:
\begin{equation*}
    \alpha_i = 0, \quad a_i = 0.25\text{ m}, \quad d_i = 0, \quad \theta_i = q_i
\end{equation*}
\begin{enumerate}
    \item Escriba la matriz de transformación genérica ${}^{i-1}T_i$.
    \item Evalúe numéricamente la matriz para una configuración articular de $q_i = 90^\circ$.
    \item Interprete físicamente la última columna de la matriz resultante.
\end{enumerate}

\section*{Ejercicio 2: Cinemática Directa Planar (3R)[cite: 1]}
Considere un manipulador planar con los siguientes parámetros geométricos y configuración articular actual:
\begin{equation*}
    a_1 = 0.30\text{ m}, \quad a_2 = 0.25\text{ m}, \quad a_3 = 0.15\text{ m}
\end{equation*}
\begin{equation*}
    q_1 = 0^\circ, \quad q_2 = 90^\circ, \quad q_3 = -90^\circ
\end{equation*}
\begin{enumerate}
    \item Construya las tres matrices de transformación ${}^{0}T_1$, ${}^{1}T_2$ y ${}^{2}T_3$ asumiendo $\alpha_i = 0$, $d_i = 0$ y $\theta_i = q_i$ para todas las juntas.
    \item Escriba la secuencia de multiplicación necesaria para obtener ${}^{0}T_3$.
    \item Calcule numéricamente la matriz resultante ${}^{0}T_3$.
    \item Extraiga el vector de posición del efector final.
    \item Realice una verificación geométrica rápida que demuestre que su posición calculada es plausible en el plano.
\end{enumerate}

\section*{Ejercicio 3: Aplicación Restringida ABB IRB 120[cite: 1]}
Usando el modelo D-H publicado por Truc \& Lam (2020), se le proporcionan las matrices individuales para la configuración $\mathbf{q} = [0, 0, 0, 0, 0, 0]^T$:

\begin{align*}
    {}^{0}T_1 &= \begin{bmatrix} 1&0&0&0 \\ 0&0&1&0 \\ 0&-1&0&0.290 \\ 0&0&0&1 \end{bmatrix}, &
    {}^{1}T_2 &= \begin{bmatrix} 0&1&0&0 \\ -1&0&0&-0.270 \\ 0&0&1&0 \\ 0&0&0&1 \end{bmatrix} \\
    {}^{2}T_3 &= \begin{bmatrix} 1&0&0&0.070 \\ 0&0&1&0 \\ 0&-1&0&0 \\ 0&0&0&1 \end{bmatrix}, &
    {}^{3}T_4 &= \begin{bmatrix} 1&0&0&0 \\ 0&0&-1&0 \\ 0&1&0&0.302 \\ 0&0&0&1 \end{bmatrix} \\
    {}^{4}T_5 &= \begin{bmatrix} 1&0&0&0 \\ 0&0&1&0 \\ 0&-1&0&0 \\ 0&0&0&1 \end{bmatrix}, &
    {}^{5}T_6 &= \begin{bmatrix} 1&0&0&0 \\ 0&1&0&0 \\ 0&0&1&0.072 \\ 0&0&0&1 \end{bmatrix}
\end{align*}

\begin{enumerate}
    \item Identifique en cuál de estas matrices se refleja explícitamente el offset D-H $\theta_2 = q_2 - \pi/2$ cuando $q_2 = 0$.
    \item Escriba la cadena de transformación completa.
    \item Calcule numéricamente ${}^{0}T_6$.
    \item Extraiga $\mathbf{p}_6$ y los tres vectores ortonormales correspondientes a los ejes de la brida.
    \item Explique por qué este resultado matemático no garantiza automáticamente coincidencia directa con el "cero de fábrica" de ABB en el controlador físico.
\end{enumerate}

\end{document}
